\section{From compatibility entropy to a regular moment}\label{tra:duality}

The next estimate allows an arbitrary joint law on compatible row and column permutations. Keeping the individual marginal entropy deficits gives a weighted inequality by entropy duality. Inverse Hausdorff--Young then converts that inequality directly into a recursion for singular moments.

Put $l=\log_2 n$ and write $T_n$ for one sweep on $n=2^l$ points. Set
\[
 Z_n(p)=\operatorname{Tr}(T_n^*T_n)^p
 =\sum_{\lambda\vdash n}D_\lambda\operatorname{tr}(T_{n,\lambda}^*T_{n,\lambda})^p.
\]
The trace includes regular multiplicities. We use the positive-power comparison and inverse Fourier bound from Lemmas~\ref{lem:positive-power-comparison} and~\ref{lem:inverse-fourier-bound}, with their unnormalized regular traces.
\label{tra:duality:matrix-inequalities}
\subsection{A sparse path second moment}
\begin{lemma}\label{tra:duality:sparse}
\label{tra:duality:sparse-statement}
There is an absolute \(b>0\) such that for all sufficiently large \(n\), on all \((n-h,\ldots)\) representations with
\(1\le h\le n^{.60}\),
\begin{equation}
                             \|T_{n,\lambda}\|\le n^{-bh}.       \label{tra:duality:eq1}
\end{equation}
\end{lemma}
\begin{proof}\label{tra:duality:sparse-proof}
\label{tra:duality:sparse-second-moment}
To see this, use the permutation representation on ordered distinct \(h\)-tuples. Every irrep in \eqref{tra:duality:eq1} occurs here and not on ordered \((h-1)\)-tuples, by Young's rule. Thus we may bound the sweep on functions in the tuple representation orthogonal to functions of any proper subset of the coordinates.

Take distinct input and output tuples \(x,y\). The paths between corresponding positions in a sweep are determined (each dimension is used just once). A path at each stage uses one of the pair switches. Say two paths touch if they use the same switch at some stage. For \(I\subseteq [h]\) let
\[
   D_I=n^{|I|}\mathbb P\{\text{sweep maps }x_i\text{ to }y_i,\ i\in I\}.
\]
This is zero if the paths call for incompatible routes. On compatible paths it is \(2^{e_I}\), where \(e_I\) counts switches used by two of the paths in \(I\): we save one bit at each such switch. In particular if one path touches none of the others in the full list, including or omitting that index has no effect on \(D_I\).

The transition density relative to uniform on tuples is \((n)_h D_{[h]}/n^h\), where \((n)_h=n(n-1)\cdots(n-h+1)\). For an operator bound on the indicated irreps, we can replace \(D_{[h]}\) here by
\[
                F=\sum_{I\subseteq[h]}(-1)^{h-|I|}D_I.
\]
All terms changed act by zero on the irreps, being kernels depending on fewer tuple coordinates. And \(F=0\) unless in the touching graph on \([h]\) there are no isolated vertices.

We estimate the second moment of $F$ over uniform distinct input and output
tuples. Let $E$ be the event that both tuples are distinct and their full
touching graph has no isolated vertices. By Cauchy--Schwarz, it suffices,
up to a factor exponential in $h$, to bound
$\mathbb E_{\rm iid}[D_I^2\mathbf1_E]$ for every $I\subseteq[h]$.
Here the expectation first samples all input and output positions
independently and uniformly. Passing from this law to uniform distinct
tuples costs at most $(n^h/(n)_h)^2\le C^h$.

One factor $D_I$ can be incorporated into the sampling law:
\[
 \mathbb E_{\rm iid}[D_I^2\mathbf1_E]
 =\mathbb E_{\mu_I}[D_I\mathbf1_E].
\]
Under $\mathbb E_{\mu_I}$, sample a sweep and all input positions independently;
for $i\in I$ take the sweep's output, and for $i\notin I$ sample an
independent uniform output. This identity uses the probability definition
of $D_I$ on all endpoint lists; its formula $2^{e_I}$ will be used only on
$E$. Conditional on the sampled sweep, the individual paths are
independent. A path indexed by $I$ is uniform among the $n$ realized deck
paths, and at most two of these occupy a given switch at a given stage.
For a path outside $I$, its independently sampled endpoints likewise make
its switch uniform among the $n/2$ switches at each stage. In either case,
its probability of touching a fixed path is at most $2l/n$. Integrating
the children of a specified forest after their parents therefore bounds
its contact probability by $(2l/n)^{\#\mathrm{edges}}$.

The remaining task is to control the second density factor without paying for every possible encounter. There is a deterministic bound on the other factor \(D_I\) still left from the square. In any list of \(s\) distinct realized deck paths the number of their mutual switches is at most \((s/2)\log_2 s\). This is the butterfly contact lemma of~\cite[Lemma~3.2]{OpenAIRouting2026}; its endpoint-counted version is $s\log_2s$. Here each shared switch saves one coin, so we use the half-sized, shared-switch count. The following entropy argument also proves exactly the local form we need. To see the deterministic bound, let a walker start uniformly on one of
the $s$ selected paths. It follows that path except at a switch shared by
two selected paths, where it chooses either continuation fairly. The
walker remains uniform among the $s$ current positions. If the paths
share $e$ switches, the expected number of fair choices is therefore
$2e/s$. Given the starting position, the final position determines the
whole butterfly route. The chain rule for entropy consequently gives
$2e/s\le\log_2 s$, proving the claimed bound. Thus if the components of the full touching graph have sizes \(s_1,\ldots,s_c\), \(D_I\le\prod_j s_j^{s_j/2}\).

Let $\Gamma$ be the event that the touching graph has no isolated
vertices. The preceding Cauchy--Schwarz and change-of-law estimates give
\[
 \mathbb E_{\rm distinct}|F|^2
 \le C_0^h\max_{I\subseteq[h]}
       \mathbb E_{\mu_I}[D_I\mathbf1_\Gamma\mathbf1_{\rm distinct}],
\]
where $\mu_I$ is the sampling law obtained by weighting once by $D_I$.
We now sum the forest witnesses while keeping the remaining density
factor. For a realized graph on $\Gamma$, choose a rooted spanning
tree in every component. If there are $c$ components, their sizes
$s_1,\ldots,s_c$ satisfy $s_j\ge2$ and $\sum_js_j=h$.
There are at most $2^h h^{h-c}$ rooted forests with $c$ components:
choose their roots, then give each nonroot a parent. For each such
forest the probability of its edges is at most $(2l/n)^{h-c}$,
and the density factor on the event that these are the actual
components is at most $\prod_js_j^{s_j/2}$. After extracting this deterministic bound, discard the exact-component
and distinctness restrictions only in the forest probability. Consequently
\[
 \mathbb E_{\mu_I}[D_I\mathbf1_\Gamma\mathbf1_{\rm distinct}]
 \le 2^h\sum_{c=1}^{\lfloor h/2\rfloor}
    \max_{\substack{s_1+\cdots+s_c=h\\s_j\ge2}}
       \prod_{j=1}^c
       \left(\frac{2lh}{n}\right)^{s_j-1}s_j^{s_j/2}.
\]
This bound is uniform in $I$. It also shows explicitly why the
extra density factor does not consume the forest probability.
For sufficiently large $n$, $2l\le n^{.01}$; since $h\le n^{.60}$,
the factor of a component of size $s\ge8$ is at most
\[
 n^{-.39(s-1)+.30s}
   =n^{-.09s+.39}\le n^{-s/25}.
\]
For $2\le s\le7$, use instead $s^{s/2}\le n^{.01s}$;
then the same bound follows from
$-.39(s-1)+.01s\le-s/25$. Thus every product in the preceding
sum is at most $n^{-h/25}$, and
\[
 \mathbb E_{\rm distinct}|F|^2
 \le h(2C_0)^h n^{-h/25}\le n^{-h/50}
\]
after enlarging the absolute size threshold. For $h=1$ the
alternating kernel is identically zero. The tuple transition
density has the additional factor $(n)_h/n^h\le1$, so its
Hilbert--Schmidt norm on the new tuple constituents is at most
$n^{-h/100}$. This proves \eqref{tra:duality:eq1}, for example
with $b=1/100$.
\end{proof}
\subsection{The compatibility entropy inequality}
\label{tra:duality:grid-definition}
Our other estimate is a trace estimate for crossing row and column sweeps. Splitting the dimensions near the middle gives a grid of size \(A\times D\) (rows indexed up to \(A\), columns up to \(D\)), \(AD=n\), both lengths within a factor 2 of \(m=\sqrt n\). We will need to estimate the event that permutations acting within rows and within columns can be commuted with changed values. We include the entropic counting bound for this event in detail. Logs in it (and in relative entropy) are natural. For probability laws $P,Q$ on a finite set, we use
\[
 \mathcal D(P\|Q)=\sum_{x:P(x)>0}P(x)\log\frac{P(x)}{Q(x)},
\]
with value $+\infty$ if $P(x)>0=Q(x)$ for some $x$, and with $0\log0=0$.

Write
\[
 R=(S_D)^A,\qquad C=(S_A)^D
\]
for the row and column groups. For \(r\in R,\ c\in C\) use notation \(r_i(j)=k,\ c_k(i)=w\). Let \(\mathcal I\) be the event that row action followed by column action can also be written in opposite order. This is exactly that for each \(j\), the \(A\) outputs \(w\) as \(i\) varies are distinct. Indeed that is necessary, and gives the first column action in the opposite-order routing, after which the row actions are also determined. Factorizations in either given order are unique.
\begin{proposition}\label{tra:duality:entropy}
\label{tra:duality:entropy-statement}
Here is the relative entropy estimate. Use \(\epsilon=1/100\). If \(Q\) is any distribution supported on \(\mathcal I\), with marginals \(Q_{r_i},Q_{c_k}\), then (for sufficiently large \(m\))
\begin{equation}
 {\cal D}(Q\|U_{R\times C})\ \ge\
 n-O(n^{.54})
 +\left(1-\frac{L}{\log m}\right)
   \left(\sum_i{\cal D}(Q_{r_i}\|U_{S_D})
                +\sum_k {\cal D}(Q_{c_k}\|U_{S_A})\right),       \label{tra:duality:eq2}
\end{equation}
for an absolute \(L\).
\end{proposition}
\begin{proof}\label{tra:duality:entropy-proof}
\label{tra:duality:clump-entropy-proof}
First put \(D_R={\cal D}(Q_r\|U_R)\), for the joint row marginal. Call entries of the \(r\) array clumped if for their \(j,k\) the number
\(N_{jk}=\#\{i:r_i(j)=k\}\) is greater than \(m^\epsilon\), and let \(T(r)\) count clumped entries. Then
\begin{equation}
                       \mathbb E_Q T\ \le\ O( (D_R+1)/\log m).       \label{tra:duality:eq3}
\end{equation}
We verify the useful strength of this bound. Under \(U_R\),
\(\mathbb E e^{z T}\le 2\) for \(z=c_0\log m\), \(c_0>0\) sufficiently small. To bound the contribution with some clumps, union-sum over sets of distinct cells \((j,k)\) and lists of \(h_{jk}>m^\epsilon\) rows mapping accordingly, using the weight \(e^{z\sum h_{jk}}\). Any compatible prescribed positions in row permutations cost probability at most \((e/D)^{\sum h_{jk}}\), by the falling-factorial bound. So after summing over rows, the factor per cell of size \(h\) is at most \((C e^z/h)^h\), negligible to arbitrary polynomial order even after summing over sizes and choosing a cell, for \(c_0<\epsilon\). Summing products proves the exponential bound, which implies \eqref{tra:duality:eq3} by relative entropy.

Use the chain rule with \(r\) revealed first. Reveal column permutations in random order (each column \(k\) given an independent uniform priority in \([0,1]\), earlier revealed first), and within each permutation in increasing order of \(i\). Besides \(D_R\) the chain rule gives at least the sum of the column marginal entropies \emph{relative to uniform}, that is the column relative entropy deficits \({\cal D}(Q_{c_k}\|U)\), plus contributions from comparing successive conditional probabilities to those in \(Q_{c_k}\) alone. We next lower bound those contributions, averaged over the orders; explicitly we are using the decomposition
\[
 {\cal D}(Q\|U_{R\times C})
 =D_R+\sum_k{\cal D}(Q_{c_k}\|U)
   +\sum_{k,i} \mathbb E\, {\cal D}(P_{\rm past}\|P).
\]
Here \(P\) gives probabilities for \(c_k(i)\) conditional just on \(c_k(i')\) for \(i'<i\); \(P_{\rm past}\) conditions also on \(r\) and earlier columns in the random order. The expectation includes \(Q\)-data (and we can average freely over orders).

There are \(u_i=A-i+1\) unused possible values inside permutation \(c_k\). Call atoms of \(P\) heavy when \(P(w)>m^\epsilon/u_i\), and write \(t\) for their mass. We have
\begin{equation}
  \sum_{k,i}\mathbb E t\le O\left(
                   \sum_k{\cal D}(Q_{c_k}\|U)/\log m\right).       \label{tra:duality:eq4}
\end{equation}
In fact the relative entropy in each internal prediction (reference uniform on \(u_i\)) bounds the heavy mass times \((\epsilon\log m-1)\), by grouping heavy and other atoms. These entropies sum to the column deficits.

Consider a step with $t\le1/2$, and let $\mathcal L$ be the light class. Write
$P_{\rm light}=P(\,\cdot\mid\mathcal L)$ and let $A_{\rm past}$ be the set of
values not forbidden by the previously exposed columns. Given the row
array, the unique $j$ with $r_i(j)=k$ determines these forbidden values:
they are the output rows already obtained from original column $j$.
Compatibility forces $P_{\rm past}$ to be supported on $A_{\rm past}$.
Splitting relative entropy between the light and heavy classes gives
\[
 \mathcal D(P_{\rm past}\|P)
 \ge P_{\rm past}(\mathcal L)
       \mathcal D(P_{\rm past}(\,\cdot\mid\mathcal L)\|P_{\rm light})
 \ge P_{\rm past}(\mathcal L)[-\log P_{\rm light}(A_{\rm past})],
\]
with the product interpreted as zero when $P_{\rm past}(\mathcal L)=0$.
After averaging over the data, the factor $P_{\rm past}(\mathcal L)$ may be
replaced by the indicator that the actual value lies in  $\mathcal L$.

Now fix all data drawn under $Q$, with $t\le1/2$ and the actual value
light, and average only over the column priorities. The internal
predictor $P$, its light class, and $P_{\rm light}$ are fixed during
this averaging. Conditional on priority $x$ for column $k$, each output
row value from another column is forbidden with probability $x$.
The exceptions have total $P_{\rm light}$-mass $\alpha$, where
\[
 \alpha\le\min(1,2N_{jk}m^\epsilon/u_i).
\]
The expected allowed mass is $1-x+x\alpha$. Jensen's inequality for
$-\log$ and then averaging over $x$ give the lower bound
$\int_0^1-\log(1-x+x\alpha)\,dx$. This argument compares with a light-class distribution to keep the loss on discarding heavy atoms proportional to their actual mass. In particular dropping \(t>1/2\) cases and actual-heavy cases loses in expectation at most three times the expected heavy mass from the potential one nat per entry, using the marginal prediction law \(P\). Dropping clumped entries loses at most \(\mathbb E T\), since as \(i,k\) vary the corresponding \(i,j\) go once through the row array.

For all remaining entries the integral bound falls short of 1 by at most \(O(a\log(e/a))\) where \(a=\min(1,2m^{2\epsilon}/u_i)\), by direct integration. Their total loss on this account is at most \(O(D m^{2\epsilon}\log^2 m)=O(n^{.54})\). Consequently the sum of the entropy comparisons is at least
\[
       n-O(n^{.54})
       -O\left((1+D_R+\sum_k{\cal D}(Q_{c_k}\|U))/\log m\right).
\]
Together with \(D_R\ge\sum_i{\cal D}(Q_{r_i}\|U)\) this proves \eqref{tra:duality:eq2}.
\end{proof}
\begin{proof}[Proof of Theorem~\ref{tra:duality:weighted}]\label{tra:duality:weighted-proof}
\label{tra:duality:entropy-duality-proof}
This is the finite-space entropy--Brascamp--Lieb duality of Carlen and Cordero-Erausquin~\cite[Theorem~2.1]{CarlenCorderoErausquin2009}. We give its specialization here. Suppose first that all weights are positive, and put
\[
 V(r,c)=\sum_i\log w_i(r_i)+\sum_k\log v_k(c_k).
\]
The finite Gibbs variational formula, with the event incorporated into the reference measure, is
\[
 \log\mathbb E_U[\mathbf1_{\mathcal I}e^V]
 =\sup_{Q:\,Q(\mathcal I)=1}\{\mathbb E_QV-\mathcal D(Q\|U)\}.
\]
Insert \eqref{tra:duality:eq2}. Each row marginal then contributes at most
\[
 \mathbb E_{Q_{r_i}}\log w_i
       -\theta\mathcal D(Q_{r_i}\|U_{S_D})
 \le\theta\log\mathbb E_{U_{S_D}}w_i^{1/\theta},
\]
and the same inequality holds for each column marginal. Dropping the requirement that those marginals arise from a common compatible law only increases the supremum. Hence
\[
 \log\mathbb E_U[\mathbf1_{\mathcal I}e^V]
 \le -n+O(n^{.54})
       +\theta\sum_i\log\mathbb Ew_i^{1/\theta}
       +\theta\sum_k\log\mathbb Ev_k^{1/\theta}.
\]
Stirling's bound gives
$\log(n!/(|R||C|))\le n+O(m\log n)$.
It cancels the leading $-n$, and $m\log n=O(n^{.54})$.
Exponentiating proves the claim for positive weights. Replace each zero weight by a positive number tending to zero to obtain the general case, since all spaces are finite. The error is independent of the weights.
\end{proof}
\subsection{Positive coefficients and compatibility}

We isolate the algebraic passage from line operators to a compatibility
probability. It applies to every rectangle; balance will enter only
through the entropy estimate used afterward. The two quantities attached
to a positive line operator $Z_v$ have different roles: the squared coefficient
density has total mass $\operatorname{Tr}Z_v^2$, while the regular rank
bounds its normalized density.

\begin{lemma}[Positive line coefficients]\label{lem:coefficient-compatibility}
Let $R,C\le S_n$ be the row and column groups of a rectangle.
For each line $v$, let $Z_v\ne0$ be a positive operator in that
line's group algebra, acting on its regular representation. Write
$z_v$ for its coefficient density relative to uniform measure, and put
\[
 w_v=\operatorname{Tr}_{\rm reg}Z_v^2,\qquad
 R_v=\operatorname{rank}_{\rm reg}Z_v,\qquad
 \rho_v=|z_v|^2/w_v.
\]
Then $\rho_v$ is a probability density and $\|\rho_v\|_\infty\le R_v$.
Let $F$ and $G$ be the tensor products of the row and column operators,
respectively, acting on the full regular representation of $S_n$.
For the compatibility event $\mathcal I=\{(r,c):cr\in RC\}$,
\begin{equation}\label{e:positive-coefficient-probability}
 \operatorname{Tr}_{\rm reg}(FGFG)
 \le\frac{n!}{|R||C|}\left(\prod_vw_v\right)
       \mathbb P_{\otimes_v\rho_v}(\mathcal I).
\end{equation}
For projections, $w_v=R_v$. A projection of carrier rank $r$ in an
irreducible of dimension $D$ has regular rank $R_v=Dr$.
\end{lemma}
\begin{proof}
For a line group $H_v$, the coefficient identity and Parseval give
\[
 z_v(g)=\operatorname{Tr}(\mathcal L(g^{-1})Z_v),\qquad
 \mathbb E_{H_v}|z_v|^2=\operatorname{Tr}Z_v^2=w_v.
\]
Since $Z_v$ is positive and $\mathcal L(g^{-1})$ is unitary,
\[
 |z_v(g)|\le\operatorname{Tr}Z_v,
 \qquad(\operatorname{Tr}Z_v)^2\le R_v\operatorname{Tr}Z_v^2.
\]
These prove the assertions about $\rho_v$.
Write $f(r)=\prod_i z_i(r_i)$ and $g(c)=\prod_k z_k(c_k)$ for
the coefficient densities of $F,G$. Their coefficients in the group
algebra are $f/|R|$ and $g/|C|$.
The trace picks out words whose permutations cancel. More precisely,
\begin{equation}\label{e:four-factor-identity}
 \Tr(FGFG)=\frac{n!}{|R|^2|C|^2}
 \sum_{\substack{r_1c_1r_2c_2=e\\r_1,r_2\in R,\ c_1,c_2\in C}}
       f(r_1)g(c_1)f(r_2)g(c_2).
\end{equation}
The factor $n!$ is the dimension of the full regular representation. Each subgroup coefficient contributes its density divided by its subgroup order.

Let $\mathcal J=\{(r,c):rc\in CR\}$. Since $R\cap C=\{e\}$, every $(r,c)\in\mathcal J$ has a unique pair $(r',c')\in R\times C$ satisfying $rc\,r'c'=e$. The map $\Psi(r,c)=(r',c')$ is an involution of $\mathcal J$: the same equality also gives $r'c'rc=e$. Cauchy--Schwarz therefore gives
\[
 \sum_{(r,c)\in\mathcal J}
   |f(r)g(c)f(r')g(c')|
 \le \sum_{(r,c)\in\mathcal J}|f(r)|^2|g(c)|^2.
\]
Finally, $(r,c)\mapsto(r^{-1},c^{-1})$ maps $\mathcal J$ bijectively to the compatibility event $\mathcal I=\{(r,c):cr\in RC\}$. Self-adjointness gives $f(r^{-1})=\overline{f(r)}$ and $g(c^{-1})=\overline{g(c)}$, so this inversion leaves the last weight unchanged. Thus
\begin{equation}\label{e:four-factor}
 \Tr(FGFG)\le\frac{n!}{|R||C|}\,
       \mathbb E_{R\times C}
          [\mathbf1_{\mathcal I}|f(r)|^2|g(c)|^2].
\end{equation}
Dividing each squared line coefficient by its mass $w_v$ gives
\eqref{e:positive-coefficient-probability}.
\end{proof}

\subsection{Trace recursion and completion}
\begin{proposition}\label{tra:duality:recursion}
\label{tra:duality:trace-recursion-statement}
\textbf{Grid trace estimate.} Suppose at the two child sizes \(A,D\) we have \(Z_{A}(p), Z_{D}(p)\le 2\), \(p\ge 4\). Splitting a sweep into independent sweeps within rows, then independent sweeps within columns, we have
\begin{equation}
        Z_n(t)\ \le\ \exp(O(n^{.54})),\qquad
                  t=p(1+O(1/\log m)),                             \label{tra:duality:eq6}
\end{equation}
where \(t\) can be taken as \(p(1+L'/\log m)\) for a suitable absolute \(L'\), for sufficiently large \(m\).
\end{proposition}
\begin{proof}\label{tra:duality:recursion-proof}
\label{tra:duality:hausdorff-young-recursion}
Write $T_n=K_CK_R$ for the chronological row--column factorization,
and set $X=K_RK_R^*$, $Y=K_C^*K_C$. The positive-power comparison
\eqref{eq:fourth-positive-trace}, with $p=t$, gives
\[
 Z_n(t)\le\operatorname{Tr}(X^{t/2}Y^{t/2}X^{t/2}Y^{t/2}).
\]
The line factors of $X^{t/2}$ and $Y^{t/2}$ have coefficient
densities $f_i$ and $g_k$. Lemma~\ref{lem:coefficient-compatibility},
in its unnormalized form \eqref{e:four-factor}, now bounds the right
side by the weighted compatibility expectation with line weights
$|f_i|^2,|g_k|^2$.
Thus we may use \eqref{tra:duality:eq5} and need only bound norms with exponent \(\nu=2/\theta\) at each fiber. By Lemma~\ref{lem:inverse-fourier-bound},
\[
 \|f_i\|_\nu
 \le \left[ Z_D\left(\tfrac{t}{2}\,\tfrac{\nu}{\nu-1}\right)
                                      \right]^{1-1/\nu}
\]
(and likewise using \(Z_A\) for \(g_k\)). The child exponent is exactly $t/(2-\theta)$. Thus $t=p(2-\theta)=p(1+L/\log m)$ makes it $p$; taking a larger absolute $L'$ also works. The norms are bounded by hypothesis, since the unpowered matrices are contractions. The product of bounds takes only \(\exp(O(m))\), establishing \eqref{tra:duality:eq6}.
\end{proof}
\begin{proof}[Proof of Theorem~\ref{tra:duality:moment}]\label{tra:duality:moment-proof}
\label{tra:duality:moment-induction-and-completion}
The trace bound is not yet small. We now split it into small diagram levels, already controlled by the forest argument, and large-dimensional blocks, where regular multiplicity amplifies any increase in the moment exponent. To justify treating all remaining representations as large after \eqref{tra:duality:eq1}, here are the dimension details. Those with first column of length greater than \(n/2\) have zero sweep block: a single parallel dimension projects onto invariants of \((S_2)^{n/2}\), which by Young's rule can occur only if \(\lambda\) dominates \((2,\ldots,2)\), thus has at most \(n/2\) rows. Among the remaining shapes, except \(\lambda=(n-h,\ldots)\) for \(0\le h\le n^{.60}\), we can use
\begin{equation}
                           D_\lambda\ \ge\ \exp(n^{.58})            \label{tra:duality:eq8}
\end{equation}
for sufficiently large \(n\). In fact if a first row or column has length at least \(n/10\), it also has at least \(s=\lfloor n^{.59}\rfloor\) boxes off it. Transposing if needed, take the full long row and a subdiagram with \(s\) boxes below. Filling the first \(s\) boxes in the long row first, we may interleave its remaining boxes and a standard ordering below, giving \eqref{tra:duality:eq8} by tableau count and branching. If neither is long, all hook lengths are bounded by \(2n/10\), giving even exponential growth in \(n\). On the other hand at levels \(h\le n^{.60}\) we may use \(D_\lambda\le n^{2h}\), immediately by branching or the ordered tuple representation.

For large \(n\) let the baseline \(p\) in \eqref{tra:duality:eq6} be at least a sufficiently large absolute constant, so that by \eqref{tra:duality:eq1} the contribution to \(Z_n(p)\) from levels 1 through \(n^{.60}\) is \(o(1)\): for each partition at level \(h\) it is at most
\(n^{4h-2pbh}\), and there are at most \(2^h\) such partitions. For the shapes in \eqref{tra:duality:eq8}, \eqref{tra:duality:eq6} forces every squared singular value there to obey
\[
                      (\text{squared singular value})^t
                         \le \exp(-\tfrac12 n^{.58})
\]
for sufficiently large \(n\), by multiplicity in the regular representation. Put $\delta=1/\log m$, and list the squared singular values in
block $\lambda$ as $q_{\lambda,j}$, $1\le j\le D_\lambda$,
including repetitions. Their total contribution over the
large-dimensional blocks at exponent $t(1+\delta)$ is bounded by
\[
 \begin{aligned}
 \sum_{\lambda\text{ large}}D_\lambda
       \sum_{j=1}^{D_\lambda}q_{\lambda,j}^{t(1+\delta)}
 &\le Z_n(t)\left(\max_{\lambda\text{ large},j}q_{\lambda,j}^t\right)^\delta\\
 &\le\exp\!\left(Cn^{.54}-\frac{n^{.58}}{2\log m}\right)=o(1).
 \end{aligned}
\]
The exponent increase therefore makes both the sparse and the
large-dimensional contributions small. We have therefore proved for sufficiently large \(n\): if the two child bounds \(Z(p)\le 2\) hold at parameter \(p\) at least the baseline constant, we get \eqref{tra:duality:eq7}'s bound at \(n\) at parameter
\[
                              p(1+L''/\log n)
\]
for an absolute \(L''\). The sufficiently-large threshold need not depend on \(p\).

For the remaining bounded sizes we can take a fixed baseline parameter large enough to get \eqref{tra:duality:eq7}; indeed no nonconstant part has norm 1, since the sweep composes orthogonal projections and their common invariants are constants (cube swaps generate all permutations). Now the parameter increases over the dimension splitting recursion stay absolutely bounded: child log sizes are about half the parent log size, so the product of \(1+L''/\log n\) over internal sizes above the threshold along any path is bounded. By monotonicity this proves \eqref{tra:duality:eq7}.

Finally take a fixed integer \(M\ge p_*\). After \(M\) sweeps the chi-squared distance to uniform is at most \(Z_n(M)-1\): indeed it is the squared Hilbert-Schmidt norm of the regular convolution matrix minus constants after taking the sweeps (all starting rows have the same density norm), and Schatten Hölder bounds this on the nonconstant space by the indicated singular moment minus the constant contribution. Thus total variation is at most $\tfrac12\sqrt{1/16}=1/8$, giving the asserted order upper bound.
\end{proof}
